% GENERATED FILE. Edit canonical lessons, then run npm run generate:books.
% canonical-source-hash: fnv1a64:413dbdaa43bd8c81
% canonical-lessons: TE-C144-vandu, TE-C144-koyu, TE-C144-pagalagottu, TE-C144-pampu, TE-C144-vetuku

\chapter{To cook, To cut, To break, To send, To look for}
\label{ch:te-to-cook-to-cut-to-break-to-send-to-look-for}
\hlchaptermodality{144}

\begin{chapteropening}
\textbf{By the end of this chapter:} \emph{I can say to cook, to cut, to break, to send and to look for.}

Complete the last lesson of chapter 144: I can say to cook, to cut, to break, to send and to look for.
\end{chapteropening}

\section[\texorpdfstring{\te{వండు} (vaṇḍu)}{వండు (vaṇḍu)}]{\te{వండు} (vaṇḍu) — to cook}

\label{lesson:TE-C144-vandu}



\begin{quote}
Before the new one: say the Telugu for to sleep, then the Telugu for to get up.
\end{quote}

\subsection*{You'll want to know: \te{వండు}}
\textbf{\te{వండు}} — \emph{vaṇḍu} — \textquotedblleft{}to cook\textquotedblright{}.

\begin{grammarlens}[title={What you've built}]
One more word for this chapter.
\end{grammarlens}

\subsection*{Your turn}
\begin{itemize}
\raggedright
  \item \emph{Say it:} \emph{vaṇḍu}
  \item \emph{Say it:} \emph{vaṇḍu}, once more
  \item \emph{Say it:} say \emph{vaṇḍu}
  \item \emph{Recall:} say the Telugu for to sleep, then the Telugu for to get up, then say \emph{vaṇḍu} again
\end{itemize}

\begin{tcolorbox}[breakable,colback=teal!4,colframe=teal!35!black,title={Before you move on}]
What does \te{వండు} mean? (\textquotedblleft{}To cook\textquotedblright{}.) Say it once more.
\end{tcolorbox}
\section[\texorpdfstring{\te{కోయు} (kōyu)}{కోయు (kōyu)}]{\te{కోయు} (kōyu) — to cut}

\label{lesson:TE-C144-koyu}



\begin{quote}
Before the new one: say the Telugu for to get up, then the Telugu for to cook.
\end{quote}

\subsection*{You'll want to know: \te{కోయు}}
\textbf{\te{కోయు}} — \emph{kōyu} — \textquotedblleft{}to cut\textquotedblright{}.

\begin{grammarlens}[title={What you've built}]
One more word for this chapter.
\end{grammarlens}

\subsection*{Your turn}
\begin{itemize}
\raggedright
  \item \emph{Say it:} \emph{kōyu}
  \item \emph{Say it:} \emph{kōyu}, once more
  \item \emph{Say it:} say \emph{kōyu}
  \item \emph{Recall:} say the Telugu for to get up, then the Telugu for to cook, then say \emph{kōyu} again
\end{itemize}

\begin{tcolorbox}[breakable,colback=teal!4,colframe=teal!35!black,title={Before you move on}]
What does \te{కోయు} mean? (\textquotedblleft{}To cut\textquotedblright{}.) Say it once more.
\end{tcolorbox}
\section[\texorpdfstring{\te{పగలగొట్టు} (pagalagoṭṭu)}{పగలగొట్టు (pagalagoṭṭu)}]{\te{పగలగొట్టు} (pagalagoṭṭu) — to break}

\label{lesson:TE-C144-pagalagottu}



\begin{quote}
Before the new one: say the Telugu for to cook, then the Telugu for to cut.
\end{quote}

\subsection*{You'll want to know: \te{పగలగొట్టు}}
\textbf{\te{పగలగొట్టు}} — \emph{pagalagoṭṭu} — \textquotedblleft{}to break\textquotedblright{}.

\begin{grammarlens}[title={What you've built}]
One more word for this chapter.
\end{grammarlens}

\subsection*{Your turn}
\begin{itemize}
\raggedright
  \item \emph{Say it:} \emph{pagalagoṭṭu}
  \item \emph{Say it:} \emph{pagalagoṭṭu}, once more
  \item \emph{Say it:} say \emph{pagalagoṭṭu}
  \item \emph{Recall:} say the Telugu for to cook, then the Telugu for to cut, then say \emph{pagalagoṭṭu} again
\end{itemize}

\begin{tcolorbox}[breakable,colback=teal!4,colframe=teal!35!black,title={Before you move on}]
What does \te{పగలగొట్టు} mean? (\textquotedblleft{}To break\textquotedblright{}.) Say it once more.
\end{tcolorbox}
\section[\texorpdfstring{\te{పంపు} (pampu)}{పంపు (pampu)}]{\te{పంపు} (pampu) — to send}

\label{lesson:TE-C144-pampu}



\begin{quote}
Before the new one: say the Telugu for to cut, then the Telugu for to break.
\end{quote}

\subsection*{You'll want to know: \te{పంపు}}
\textbf{\te{పంపు}} — \emph{pampu} — \textquotedblleft{}to send\textquotedblright{}.

\begin{grammarlens}[title={What you've built}]
One more word for this chapter.
\end{grammarlens}

\subsection*{Your turn}
\begin{itemize}
\raggedright
  \item \emph{Say it:} \emph{pampu}
  \item \emph{Say it:} \emph{pampu}, once more
  \item \emph{Say it:} say \emph{pampu}
  \item \emph{Recall:} say the Telugu for to cut, then the Telugu for to break, then say \emph{pampu} again
\end{itemize}

\begin{tcolorbox}[breakable,colback=teal!4,colframe=teal!35!black,title={Before you move on}]
What does \te{పంపు} mean? (\textquotedblleft{}To send\textquotedblright{}.) Say it once more.
\end{tcolorbox}
\section[\texorpdfstring{\te{వెతుకు} (vetuku)}{వెతుకు (vetuku)}]{\te{వెతుకు} (vetuku) — to look for}

\label{lesson:TE-C144-vetuku}



\begin{quote}
Before the new one: say the Telugu for to break, then the Telugu for to send.
\end{quote}

\subsection*{You'll want to know: \te{వెతుకు}}
\textbf{\te{వెతుకు}} — \emph{vetuku} — \textquotedblleft{}to look for\textquotedblright{}.

\begin{grammarlens}[title={What you've built}]
One more word for this chapter.
\end{grammarlens}

\subsection*{Your turn}
\begin{itemize}
\raggedright
  \item \emph{Say it:} \emph{vetuku}
  \item \emph{Say it:} \emph{vetuku}, once more
  \item \emph{Say it:} say \emph{vetuku}
  \item \emph{Recall:} say the Telugu for to break, then the Telugu for to send, then say \emph{vetuku} again
\end{itemize}

\begin{tcolorbox}[breakable,colback=teal!4,colframe=teal!35!black,title={Before you move on}]
What does \te{వెతుకు} mean? (\textquotedblleft{}To look for\textquotedblright{}.) Say it once more.
\end{tcolorbox}
